\section{Law of Large Numbers}

\subsection{Convergence in probability}

Throughout this section, random variables that agree almost surely are
identified.

\begin{de}[Convergence in probability]
A sequence $(X_n)$ converges to $X$ in probability, written
$X_n\xrightarrow{P}X$, if for every $\varepsilon>0$,
\[
P(|X_n-X|>\varepsilon)\longrightarrow0.
\]
\end{de}

\begin{prop}[A metric for convergence in probability]
On $L^0(\Omega,\mathcal F,P)$, define
\[
d_0(X,Y)=\mathbb E\bigl[1\wedge|X-Y|\bigr].
\]
Then $d_0$ is a metric, and
\[
d_0(X_n,X)\to0
\quad\Longleftrightarrow\quad
X_n\xrightarrow{P}X.
\]
\end{prop}

\begin{proof}
Only the triangle inequality and the convergence equivalence require
comment. Since $1\wedge(a+b)\leq(1\wedge a)+(1\wedge b)$, the triangle
inequality follows pointwise. If $d_0(X_n,X)\to0$, then for
$0<\varepsilon\leq1$,
\[
P(|X_n-X|>\varepsilon)
 \leq\varepsilon^{-1}d_0(X_n,X).
\]
Conversely, convergence in probability implies that
$1\wedge|X_n-X|\to0$ in probability. Every subsequence has an almost surely
convergent further subsequence, and dominated convergence shows that every
subsequence of $d_0(X_n,X)$ has a further subsequence tending to zero. Hence
the whole numerical sequence tends to zero.
\end{proof}

\begin{thm}[Completeness of $L^0$]
The metric space $(L^0,d_0)$ is complete.
\end{thm}

\begin{proof}
Let $(X_n)$ be Cauchy in probability. Choose a subsequence $(X_{n_k})$ such
that
\[
P(|X_{n_{k+1}}-X_{n_k}|>2^{-k})\leq2^{-k}.
\]
The first Borel--Cantelli lemma implies that, almost surely, only finitely many
events
\[
\{|X_{n_{k+1}}-X_{n_k}|>2^{-k}\}
\]
occur. Thus $(X_{n_k})$ is almost surely Cauchy and converges to a finite
random variable $X$. It follows that
$X_{n_k}\to X$ in probability. Since the original sequence is Cauchy in
probability, the whole sequence converges to $X$ in probability.
\end{proof}

\begin{prop}[Basic properties]
Let $X_n\xrightarrow{P}X$ and $Y_n\xrightarrow{P}Y$.
\begin{enumerate}[label=(\roman*)]
  \item The limit in probability is unique up to almost sure equality.
  \item $X_n+Y_n\xrightarrow{P}X+Y$ and
  $X_nY_n\xrightarrow{P}XY$.
  \item If $f:\mathbb R^d\to\mathbb R^m$ is continuous and
  $Z_n\xrightarrow{P}Z$, then $f(Z_n)\xrightarrow{P}f(Z)$.
  \item Almost sure convergence implies convergence in probability.
  \item $X_n\xrightarrow{P}X$ if and only if every subsequence has a further
  subsequence that converges to $X$ almost surely.
\end{enumerate}
\end{prop}

\begin{proof}
Uniqueness follows from
$P(|X-Y|>\varepsilon)\leq
P(|X-X_n|>\varepsilon/2)+P(|X_n-Y|>\varepsilon/2)$.
The sum follows from the same estimate. For products, first note that a
sequence converging in probability is bounded in probability, then write
$X_nY_n-XY=X_n(Y_n-Y)+Y(X_n-X)$. The continuous-mapping assertion follows by
a subsequence argument, and the last two assertions follow respectively from
dominated convergence for indicators and the Borel--Cantelli subsequence
construction used above.
\end{proof}

\subsection{Independence}

\begin{de}[Independent families]
Families of events $\mathcal C_i\subseteq\mathcal F$, $i\in I$, are
independent if, for every finite set of distinct indices
$i_1,\ldots,i_m$ and every $A_j\in\mathcal C_{i_j}$,
\[
P\!\left(\bigcap_{j=1}^{m}A_j\right)
 =\prod_{j=1}^{m}P(A_j).
\]
Random variables $(X_i)_{i\in I}$ are independent if the $\sigma$-fields
$\sigma(X_i)$ are independent.
\end{de}

\begin{de}[$\pi$-systems and Dynkin systems]
A nonempty family $\mathcal P$ is a $\pi$-system if it is closed under finite
intersections. A family $\mathcal D$ is a Dynkin system if it contains
$\Omega$, is closed under complements, and is closed under countable disjoint
unions.
\end{de}

\begin{thm}[$\pi$--$\lambda$ theorem]
If a Dynkin system $\mathcal D$ contains a $\pi$-system $\mathcal P$, then
$\sigma(\mathcal P)\subseteq\mathcal D$.
\end{thm}

\begin{cor}[Independence passes to generated $\sigma$-fields]
If $\mathcal C_1,\ldots,\mathcal C_m$ are independent $\pi$-systems, each
containing $\Omega$, then
$\sigma(\mathcal C_1),\ldots,\sigma(\mathcal C_m)$ are independent.
\end{cor}

\begin{proof}
Fix all but one coordinate and consider the class of sets in the remaining
coordinate for which the required product identity holds. It is a Dynkin
system containing the corresponding $\pi$-system. Apply the
$\pi$--$\lambda$ theorem, and repeat one coordinate at a time.
\end{proof}

\begin{lem}[Second Borel--Cantelli lemma]
If $(A_n)$ is an independent sequence of events and
$\sum_nP(A_n)=\infty$, then
\[
P(A_n\ \mathrm{i.o.})=1.
\]
\end{lem}

\begin{proof}
For $m\leq N$, independence and $1-x\leq e^{-x}$ give
\[
P\!\left(\bigcap_{n=m}^{N}A_n^c\right)
 =\prod_{n=m}^{N}(1-P(A_n))
 \leq\exp\!\left(-\sum_{n=m}^{N}P(A_n)\right).
\]
Let $N\to\infty$. The right-hand side tends to zero, so
$P(\bigcup_{n\geq m}A_n)=1$ for every $m$. Intersect over $m$.
\end{proof}

\begin{prop}[Measurable transformations preserve independence]
If $X_1,\ldots,X_m$ are independent and $f_i$ are measurable, then
$f_1(X_1),\ldots,f_m(X_m)$ are independent.
\end{prop}

\begin{proof}
For every Borel set $B$, the event
$\{f_i(X_i)\in B\}$ belongs to $\sigma(X_i)$.
\end{proof}

\begin{thm}[Product-law characterization]
Random variables $X_i:(\Omega,\mathcal F,P)\to(S_i,\mathcal S_i)$,
$i=1,\ldots,m$, are independent if and only if
\[
\mathcal L(X_1,\ldots,X_m)
 =\mathcal L(X_1)\otimes\cdots\otimes\mathcal L(X_m).
\]
Equivalently, for all bounded measurable $f_i$,
\[
\mathbb E\!\left[\prod_{i=1}^{m}f_i(X_i)\right]
 =\prod_{i=1}^{m}\mathbb E[f_i(X_i)].
\]
The same identity holds for nonnegative measurable functions, with values in
$[0,\infty]$.
\end{thm}

\begin{proof}
Independence gives the product identity on measurable rectangles, which form
a generating $\pi$-system. Uniqueness of product measures gives the law
identity. The integral factorization follows first for indicators, then for
simple functions, and finally by bounded or monotone convergence. Conversely,
apply the functional identity to indicators.
\end{proof}

\begin{prop}[Equivalent tests for real-valued independence]
For real random variables $X_1,\ldots,X_m$, the following are equivalent:
\begin{enumerate}[label=(\roman*)]
  \item the variables are independent;
  \item the bounded-measurable factorization above holds;
  \item it holds for all bounded continuous $f_i$;
  \item it holds for all $f_i\in C_c(\mathbb R)$ together with constants;
  \item for all $x_1,\ldots,x_m\in\mathbb R$,
  \[
  P(X_1\leq x_1,\ldots,X_m\leq x_m)
   =\prod_{i=1}^{m}P(X_i\leq x_i).
  \]
\end{enumerate}
\end{prop}

\begin{proof}
The forward implications are immediate. The classes of test functions in
(iii) and (iv) determine finite Borel measures by the functional
monotone-class theorem. The lower orthants in (v) form a $\pi$-system that
generates $\mathcal B(\mathbb R^m)$, so equality there determines the joint
law.
\end{proof}

\begin{lem}[Grouping independent variables]
Let $(X_i)_{i\in I}$ be independent, and let $(I_r)_{r\in R}$ be pairwise
disjoint subsets of $I$. Then the $\sigma$-fields
\[
\mathcal G_r=\sigma(X_i:i\in I_r),\qquad r\in R,
\]
are independent. Consequently, measurable functions of disjoint groups of
independent variables are independent.
\end{lem}

\begin{proof}
For each $r$, finite intersections of events from
$\{\sigma(X_i):i\in I_r\}$ form a $\pi$-system generating $\mathcal G_r$.
These $\pi$-systems are independent by the definition of independence of the
whole family. Apply the preceding $\pi$--$\lambda$ argument.
\end{proof}

\begin{thm}[Moment characterization for bounded variables]
Let $X_1,\ldots,X_m$ be bounded real random variables. They are independent
if and only if, for every $k_1,\ldots,k_m\in\mathbb N_0$,
\[
\mathbb E[X_1^{k_1}\cdots X_m^{k_m}]
 =\prod_{i=1}^{m}\mathbb E[X_i^{k_i}].
\]
\end{thm}

\begin{proof}
Necessity follows from the product-law characterization. Conversely, choose
a compact interval containing the ranges of all variables. By linearity, the
factorization holds for polynomials in each coordinate. The Weierstrass
approximation theorem extends it uniformly to continuous functions on that
interval. The bounded-continuous characterization of independence then
applies.
\end{proof}

\subsection{Weak law of large numbers}

\begin{de}[Uncorrelated variables]
Square-integrable random variables $(X_i)$ are uncorrelated if
$\operatorname{Cov}(X_i,X_j)=0$ whenever $i\neq j$.
\end{de}

\begin{prop}
If $X_1,\ldots,X_n$ are pairwise uncorrelated, then
\[
\operatorname{Var}\!\left(\sum_{i=1}^{n}X_i\right)
 =\sum_{i=1}^{n}\operatorname{Var}(X_i).
\]
\end{prop}

\begin{proof}
Expand the square of the centered sum. All cross-covariance terms vanish.
\end{proof}

\begin{thm}[$L^2$ weak law]
Let $(X_n)$ be pairwise uncorrelated with common mean $\mu$ and
$\sup_n\operatorname{Var}(X_n)\leq C<\infty$. If
$S_n=\sum_{k=1}^{n}X_k$, then
\[
\frac{S_n}{n}\longrightarrow\mu
\quad\text{in }L^2\text{ and in probability}.
\]
\end{thm}

\begin{proof}
The mean of $S_n/n$ is $\mu$, and
\[
\mathbb E\left|\frac{S_n}{n}-\mu\right|^2
 =\frac1{n^2}\sum_{k=1}^{n}\operatorname{Var}(X_k)
 \leq\frac Cn\to0.
\]
\end{proof}

\begin{prop}[Chebyshev form]
If $S_n$ is square integrable, $a_n=\mathbb E S_n$, and
$\operatorname{Var}(S_n)/b_n^2\to0$, then
\[
\frac{S_n-a_n}{b_n}\xrightarrow{P}0.
\]
\end{prop}

\begin{thm}[Weak law for triangular arrays]
For each $n$, let $X_{n,1},\ldots,X_{n,n}$ be independent. Let $b_n\to\infty$
and set
\[
\overline X_{n,k}=X_{n,k}\mathbf1_{\{|X_{n,k}|\leq b_n\}},
\qquad
a_n=\sum_{k=1}^{n}\mathbb E\overline X_{n,k}.
\]
Assume
\[
\sum_{k=1}^{n}P(|X_{n,k}|>b_n)\to0
\]
and
\[
\frac1{b_n^2}\sum_{k=1}^{n}\mathbb E[\overline X_{n,k}^{\,2}]\to0.
\]
Then, with $S_n=\sum_{k=1}^{n}X_{n,k}$,
\[
\frac{S_n-a_n}{b_n}\xrightarrow{P}0.
\]
\end{thm}

\begin{proof}
Let $\overline S_n=\sum_k\overline X_{n,k}$. The union bound gives
\[
P(S_n\neq\overline S_n)
 \leq\sum_{k=1}^{n}P(|X_{n,k}|>b_n)\to0.
\]
Independence and Chebyshev's inequality give
\[
P(|\overline S_n-a_n|>\varepsilon b_n)
 \leq\frac1{\varepsilon^2b_n^2}
      \sum_{k=1}^{n}\operatorname{Var}(\overline X_{n,k})
 \leq\frac1{\varepsilon^2b_n^2}
      \sum_{k=1}^{n}\mathbb E[\overline X_{n,k}^{\,2}]
 \to0.
\]
Combine the two estimates.
\end{proof}

\begin{thm}[Feller-type weak law]
Let $(X_n)$ be i.i.d. and suppose
\[
xP(|X_1|>x)\longrightarrow0.
\]
Put $S_n=\sum_{k=1}^{n}X_k$ and
\[
\mu_n=\mathbb E[X_1\mathbf1_{\{|X_1|\leq n\}}].
\]
Then
\[
\frac{S_n}{n}-\mu_n\xrightarrow{P}0.
\]
\end{thm}

\begin{proof}
Apply the triangular-array theorem with $X_{n,k}=X_k$ and $b_n=n$. The first
condition is $nP(|X_1|>n)\to0$. For the second, use the layer-cake formula:
\[
\frac1n\mathbb E[X_1^2\mathbf1_{\{|X_1|\leq n\}}]
 \leq\frac2n\int_0^n yP(|X_1|>y)\,\dd y.
\]
The integrand $g(y)=2yP(|X_1|>y)$ tends to zero and is bounded. Its Ces\`aro
average $n^{-1}\int_0^ng(y)\,\dd y$ therefore tends to zero.
\end{proof}

\begin{cor}[Khinchin's weak law]
If $(X_n)$ is i.i.d. and $\mathbb E|X_1|<\infty$, then
\[
\frac{S_n}{n}\xrightarrow{P}\mathbb E X_1.
\]
\end{cor}

\begin{proof}
We have
$xP(|X_1|>x)\leq\mathbb E[|X_1|\mathbf1_{\{|X_1|>x\}}]\to0$, and dominated
convergence gives
$\mu_n\to\mathbb E X_1$. Apply the preceding theorem.
\end{proof}

\subsection{Independent random variables and tail events}

\begin{de}[Tail $\sigma$-field]
For a sequence $(X_n)$, the tail $\sigma$-field is
\[
\mathcal T=\bigcap_{m=1}^{\infty}\sigma(X_m,X_{m+1},\ldots).
\]
An event in $\mathcal T$ is unaffected by changing finitely many coordinates.
\end{de}

\begin{eg}
If $b_n\to\infty$, then each event
\[
\left\{\lim_{n\to\infty}\frac{S_n}{b_n}=c\right\}
\]
is a tail event: changing finitely many summands changes $S_n/b_n$ by a
quantity tending to zero. By contrast, an event such as
$\{\limsup_nS_n>0\}$ need not be a tail event, because a finite change can
shift every later partial sum by a nonvanishing constant.
\end{eg}

\begin{thm}[Kolmogorov's zero--one law]
If $(X_n)$ is independent, then every $A\in\mathcal T$ satisfies
$P(A)\in\{0,1\}$.
\end{thm}

\begin{proof}
For each $m$, $A$ is measurable with respect to
$\sigma(X_m,X_{m+1},\ldots)$ and is therefore independent of
$\sigma(X_1,\ldots,X_{m-1})$. Hence $A$ is independent of the algebra
$\bigcup_m\sigma(X_1,\ldots,X_m)$, and a $\pi$--$\lambda$ argument shows that
$A$ is independent of $\sigma(X_1,X_2,\ldots)$. In particular, $A$ is
independent of itself, so $P(A)=P(A)^2$.
\end{proof}

\begin{eg}[Borel--Cantelli zero--one law]
For independent events $(A_n)$,
\[
P(A_n\ \mathrm{i.o.})=
\begin{cases}
0,&\sum_nP(A_n)<\infty,\\
1,&\sum_nP(A_n)=\infty.
\end{cases}
\]
The event $\{A_n\ \mathrm{i.o.}\}$ is a tail event, and the two cases are the
first and second Borel--Cantelli lemmas.
\end{eg}

\begin{thm}[Approximation by finitely many coordinates]
Let $Y$ be bounded and $\sigma(X_1,X_2,\ldots)$-measurable. For every
$\varepsilon>0$, there are $N\in\mathbb N$ and a bounded Borel function
$F:\mathbb R^N\to\mathbb R$ such that
\[
\mathbb E|Y-F(X_1,\ldots,X_N)|<\varepsilon.
\]
\end{thm}

\begin{proof}
The union $\bigcup_N\sigma(X_1,\ldots,X_N)$ is an algebra generating
$\sigma(X_1,X_2,\ldots)$. Indicators of sets from this algebra are dense in
$L^1$ among indicators of the generated $\sigma$-field, and hence their
simple-function span is dense in $L^1$. Choose a simple approximation
measurable with respect to one finite-coordinate $\sigma$-field, taking $N$
to be the largest coordinate used. The Doob--Dynkin factorization gives the
Borel function $F$.
\end{proof}

\begin{cor}[Tail-measurable variables are constant]
If $(X_n)$ is independent and an extended-real random variable $Y$ is
$\mathcal T$-measurable, then $Y=c$ almost surely for some
$c\in[-\infty,\infty]$.
\end{cor}

\begin{proof}
For every $t\in\mathbb R$, the event $\{Y\leq t\}$ is a tail event, so its
probability is zero or one. Let
$c=\inf\{t\in\mathbb R:P(Y\leq t)=1\}$, allowing infinite values. Monotonicity
and right-continuity of the distribution function show that
$P(Y=c)=1$ when $c$ is finite; the cases $c=\pm\infty$ follow in the same way.
\end{proof}

\begin{de}[Convolution]
For Borel probability measures $\mu,\nu$ on $\mathbb R^d$, their convolution
is
\[
(\mu*\nu)(B)=\int_{\mathbb R^d}\mu(B-y)\,\nu(\dd y).
\]
If $X\sim\mu$ and $Y\sim\nu$ are independent, then
$X+Y\sim\mu*\nu$.
\end{de}

\begin{proof}
For Borel $B$, Tonelli's theorem and the product law give
\[
P(X+Y\in B)
 =\iint\mathbf1_B(x+y)\,\mu(\dd x)\nu(\dd y)
 =\int\mu(B-y)\,\nu(\dd y).
\]
\end{proof}

\begin{prop}[Convolution densities]
If $\mu$ and $\nu$ have Lebesgue densities $f$ and $g$, then $\mu*\nu$ has
density
\[
(f*g)(u)=\int_{\mathbb R^d}f(x)g(u-x)\,\dd x.
\]
For point masses, the corresponding formula is
\[
(\mu*\nu)(\{u\})
 =\sum_x\mu(\{x\})\nu(\{u-x\}),
\]
where only atoms contribute.
\end{prop}

\begin{proof}
Use Tonelli's theorem and the translation invariance of Lebesgue measure:
\[
(\mu*\nu)(B)
 =\iint\mathbf1_B(x+y)f(x)g(y)\,\dd x\dd y
 =\int_B\!\left(\int f(x)g(u-x)\,\dd x\right)\dd u.
\]
The atomic identity follows directly from the definition.
\end{proof}

\subsection{Strong law of large numbers}

\begin{de}[Tail equivalence]
Two sequences $(X_n)$ and $(X_n')$ on the same probability space are
\emph{tail equivalent} if
\[
\sum_{n=1}^{\infty}P(X_n\neq X_n')<\infty.
\]
Then $X_n=X_n'$ for all sufficiently large $n$, almost surely.
\end{de}

\begin{cor}
If $(X_n)$ and $(X_n')$ are tail equivalent and $b_n\to\infty$, then
\[
\frac1{b_n}\sum_{k=1}^{n}X_k
\quad\text{and}\quad
\frac1{b_n}\sum_{k=1}^{n}X_k'
\]
have the same almost sure convergence behavior and, when they converge, the
same limit.
\end{cor}

\begin{proof}
Outside a null set, the two sequences differ only in finitely many terms.
Their partial sums therefore differ by an eventually constant finite random
quantity, which divided by $b_n$ tends to zero.
\end{proof}

\begin{lem}[Integrable tails]
If $X\in L^1$ and $\varepsilon>0$, then
\[
\sum_{n=1}^{\infty}P(|X|\geq n\varepsilon)
 \leq\frac1\varepsilon\mathbb E|X|.
\]
Consequently, if $(X_n)$ is i.i.d. and integrable, then it is tail equivalent
to
$X_n'=X_n\mathbf1_{\{|X_n|\leq n\}}$.
\end{lem}

\begin{proof}
For $x\geq0$,
$\sum_{n\ge1}\mathbf1_{[n,\infty)}(x)\leq x$. Apply this to
$|X|/\varepsilon$ and take expectations. The consequence follows from the
first Borel--Cantelli lemma.
\end{proof}

\begin{thm}[Kolmogorov maximal inequality]
Let $Y_1,\ldots,Y_n$ be independent, mean-zero, square-integrable random
variables, and put $S_k=\sum_{j=1}^{k}Y_j$ and
$S_n^*=\max_{k\leq n}|S_k|$. Then, for every $\varepsilon>0$,
\[
\varepsilon^2P(S_n^*\geq\varepsilon)
 \leq\mathbb E[S_n^2\mathbf1_{\{S_n^*\geq\varepsilon\}}]
 \leq\mathbb E S_n^2.
\]
\end{thm}

\begin{proof}
Let $\tau=\inf\{k:|S_k|\geq\varepsilon\}$ and decompose the event
$\{S_n^*\geq\varepsilon\}$ into the disjoint events $\{\tau=j\}$,
$j\leq n$. Since $\{\tau=j\}\in\sigma(Y_1,\ldots,Y_j)$ and
$S_n-S_j$ is independent of this $\sigma$-field with mean zero,
\[
\mathbb E[S_n^2\mathbf1_{\{\tau=j\}}]
 =\mathbb E[S_j^2\mathbf1_{\{\tau=j\}}]
  +\mathbb E[(S_n-S_j)^2\mathbf1_{\{\tau=j\}}]
 \geq\varepsilon^2P(\tau=j).
\]
Sum over $j$.
\end{proof}

\begin{thm}[Kolmogorov convergence criterion]
Let $(Y_n)$ be independent square-integrable random variables. If
\[
\sum_{n=1}^{\infty}\operatorname{Var}(Y_n)<\infty,
\]
then
\[
\sum_{n=1}^{\infty}(Y_n-\mathbb EY_n)
\]
converges almost surely and in $L^2$. If in addition the numerical series
$\sum_n\mathbb EY_n$ converges, then $\sum_nY_n$ converges almost surely and
in $L^2$.
\end{thm}

\begin{proof}
The centered partial sums are Cauchy in $L^2$ by orthogonality. For almost
sure convergence, let $S_n=\sum_{k=1}^{n}(Y_k-\mathbb EY_k)$. Kolmogorov's
inequality, first on a finite interval and then with the endpoint tending to
infinity, gives
\[
P\!\left(\sup_{j\geq m}|S_j-S_m|>\varepsilon\right)
 \leq\frac1{\varepsilon^2}
      \sum_{j=m+1}^{\infty}\operatorname{Var}(Y_j)
 \longrightarrow0.
\]
Since
\[
\sup_{j,k\geq m}|S_j-S_k|
 \leq2\sup_{j\geq m}|S_j-S_m|,
\]
the decreasing random variables on the left converge in probability to zero.
Their pointwise limit must therefore vanish almost surely. Thus $(S_n)$ is
almost surely Cauchy. Adding a convergent
deterministic series proves the final assertion.
\end{proof}

\begin{lem}[Kronecker's lemma]
Let $b_n>0$ be nondecreasing with $b_n\to\infty$. If the series
$\sum_nx_n/b_n$ converges, then
\[
\frac1{b_n}\sum_{k=1}^{n}x_k\longrightarrow0.
\]
\end{lem}

\begin{proof}
Let $s_n=\sum_{k=1}^{n}x_k/b_k$ and $s_n\to s$. Summation by parts gives
\[
\frac1{b_n}\sum_{k=1}^{n}x_k
 =s_n-\frac1{b_n}\sum_{k=1}^{n-1}(b_{k+1}-b_k)s_k.
\]
Subtract and add $s$. The constant part tends to zero because the weights
sum to $(b_n-b_1)/b_n$, while the remainder tends to zero by splitting the
weighted sum at a fixed large index and using $s_k\to s$.
\end{proof}

\begin{thm}[Kolmogorov strong law]
Let $(X_n)$ be i.i.d. with $\mathbb E|X_1|<\infty$, and put
$S_n=\sum_{k=1}^{n}X_k$. Then
\[
\frac{S_n}{n}\longrightarrow\mathbb E X_1
\qquad\text{almost surely}.
\]
\end{thm}

\begin{proof}
Let $X_n'=X_n\mathbf1_{\{|X_n|\leq n\}}$. The integrable-tail lemma shows
that $(X_n)$ and $(X_n')$ are tail equivalent. Moreover,
\[
\sum_{n=1}^{\infty}\frac{\operatorname{Var}(X_n')}{n^2}
 \leq\mathbb E\!\left[X_1^2
     \sum_{n\geq |X_1|}\frac1{n^2}\right]
 \leq C\mathbb E|X_1|<\infty,
\]
where the summand at $X_1=0$ is interpreted as zero and the elementary bound
$\sum_{n\geq x}n^{-2}\leq C/(1+x)$ is used. The convergence criterion and
Kronecker's lemma therefore give
\[
\frac1n\sum_{k=1}^{n}(X_k'-\mathbb EX_k')\to0
\quad\text{a.s.}
\]
Dominated convergence yields $\mathbb EX_k'\to\mathbb EX_1$, so Ces\`aro's
lemma gives $n^{-1}\sum_{k\leq n}\mathbb EX_k'\to\mathbb EX_1$. Tail
equivalence completes the proof.
\end{proof}

\begin{thm}[Marcinkiewicz--Zygmund strong law]
Let $1<p<2$, let $(X_n)$ be i.i.d., and assume
$\mathbb E|X_1|^p<\infty$. Then
\[
\frac{S_n-n\mathbb EX_1}{n^{1/p}}\longrightarrow0
\qquad\text{almost surely}.
\]
\end{thm}

\begin{proof}
Set $Y_n=X_n\mathbf1_{\{|X_n|\leq n^{1/p}\}}$. Since
\[
\sum_nP(X_n\neq Y_n)
 =\sum_nP(|X_1|^p>n)\leq\mathbb E|X_1|^p,
\]
the two sequences are tail equivalent. Also,
\[
\sum_{n=1}^{\infty}\frac{\operatorname{Var}(Y_n)}{n^{2/p}}
 \leq\mathbb E\!\left[X_1^2
      \sum_{n\geq |X_1|^p}n^{-2/p}\right]
 \leq C_p\mathbb E|X_1|^p<\infty.
\]
The convergence criterion and Kronecker's lemma, with $b_n=n^{1/p}$, imply
\[
\frac1{n^{1/p}}\sum_{k=1}^{n}(Y_k-\mathbb EY_k)\to0
\quad\text{a.s.}
\]
It remains to control the means. By Tonelli,
\[
\frac1{n^{1/p}}
\sum_{k=1}^{n}\mathbb E\bigl[|X_1|
 \mathbf1_{\{|X_1|>k^{1/p}\}}\bigr]
\leq
\mathbb E\!\left[
 \frac{|X_1|\min(n,|X_1|^p)}{n^{1/p}}
\right].
\]
Split according to $|X_1|\leq n^{1/p}$ and its complement. On the first set,
the integrand is at most $|X_1|^p$ and tends pointwise to zero; on the second,
it is also at most $|X_1|^p$, and its expectation tends to zero. Dominated
convergence finishes the mean correction, and tail equivalence finishes the
proof.
\end{proof}

\begin{thm}[A weighted strong law]
Let $(X_n)$ be independent with common mean $\alpha$ and
$\sup_n\operatorname{Var}(X_n)\leq s^2$. Let $(b_n)$ be nondecreasing,
$b_n\to\infty$, and assume
\[
\sum_{n=1}^{\infty}\frac1{b_n^2}<\infty.
\]
Then
\[
\frac{\sum_{k=1}^{n}(X_k-\alpha)}{b_n}\longrightarrow0
\quad\text{almost surely and in }L^2.
\]
\end{thm}

\begin{proof}
The series
$\sum_n(X_n-\alpha)/b_n$ converges almost surely by Kolmogorov's criterion,
so Kronecker's lemma gives the almost sure conclusion. For $L^2$ convergence,
\[
\mathbb E\left|\frac{\sum_{k=1}^{n}(X_k-\alpha)}{b_n}\right|^2
 \leq s^2\frac n{b_n^2}.
\]
Because $b_n$ is nondecreasing, $a_n=b_n^{-2}$ is nonincreasing; convergence
of $\sum_na_n$ implies $na_n\to0$. Thus the displayed bound tends to zero.
\end{proof}

\subsection{Applications: renewal theory}

\begin{thm}[Elementary renewal law]
Let $(X_n)$ be i.i.d., nonnegative, integrable, and suppose
$P(X_1>0)>0$. Put $S_0=0$, $S_n=\sum_{k=1}^{n}X_k$, and
\[
N_t=\max\{n\geq0:S_n\leq t\}.
\]
Then, with $\mu=\mathbb EX_1>0$,
\[
\frac{N_t}{t}\longrightarrow\frac1\mu
\qquad\text{almost surely as }t\to\infty.
\]
\end{thm}

\begin{proof}
The strong law gives $S_n/n\to\mu>0$, so $S_n\to\infty$ and
$N_t\to\infty$ almost surely. By definition,
\[
S_{N_t}\leq t<S_{N_t+1}.
\]
Divide by $N_t$ and use the strong law along the random subsequences
$N_t$ and $N_t+1$ to obtain $t/N_t\to\mu$.
\end{proof}

\begin{lem}
If $(X_n)$ is i.i.d. and $X_1\in L^1$, then
\[
\frac{X_n}{n}\longrightarrow0
\qquad\text{almost surely}.
\]
\end{lem}

\begin{proof}
For every $\varepsilon>0$,
$\sum_nP(|X_n|\geq n\varepsilon)<\infty$ by the integrable-tail lemma. The
first Borel--Cantelli lemma gives the result first for rational
$\varepsilon>0$, and then for all $\varepsilon$.
\end{proof}

\begin{thm}[Alternating renewal reward]
Let $(X_n,Y_n)_{n\geq1}$ be i.i.d. pairs of nonnegative integrable random
variables with $P(X_1+Y_1>0)>0$. Interpret $X_n$ as an ``on'' time and $Y_n$
as the following ``off'' time. Put
\[
S_n=\sum_{k=1}^{n}(X_k+Y_k),
\qquad
N_t=\max\{n:S_n\leq t\},
\]
and
\[
T_t=\sum_{k=1}^{N_t}X_k
    +\min\{X_{N_t+1},t-S_{N_t}\}.
\]
Then
\[
\frac{T_t}{t}\longrightarrow
\frac{\mathbb EX_1}{\mathbb EX_1+\mathbb EY_1}
\qquad\text{almost surely}.
\]
\end{thm}

\begin{proof}
Let $\widetilde T_t=\sum_{k=1}^{N_t}X_k$. Then
$0\leq T_t-\widetilde T_t\leq X_{N_t+1}$. The preceding lemma and the renewal
law imply
\[
\frac{X_{N_t+1}}t
 =\frac{X_{N_t+1}}{N_t+1}\frac{N_t+1}{t}\to0.
\]
Moreover, the strong law and renewal law give
\[
\frac{\widetilde T_t}{t}
 =\frac1{N_t}\sum_{k=1}^{N_t}X_k\cdot\frac{N_t}{t}
 \longrightarrow
 \mathbb EX_1\cdot
 \frac1{\mathbb E(X_1+Y_1)}.
\]
The squeeze theorem proves the claim.
\end{proof}

\subsection{Large deviations}

\begin{de}[Exponential integrability and cumulants]
A random variable $X$ is \emph{exponentially integrable} if
$\mathbb E e^{t|X|}<\infty$ for every $t>0$. Its moment-generating function
and cumulant-generating function are
\[
M_X(t)=\mathbb E e^{tX},
\qquad
\psi_X(t)=\log M_X(t).
\]
They are finite and real analytic on $\mathbb R$, and
\[
M_X(t)=\sum_{k=0}^{\infty}\mathbb E[X^k]\frac{t^k}{k!}.
\]
The cumulants are $\kappa_k(X)=\psi_X^{(k)}(0)$; in particular,
$\kappa_1(X)=\mathbb EX$ and $\kappa_2(X)=\operatorname{Var}(X)$.
\end{de}

\begin{de}[Legendre--Fenchel transform]
For a finite convex function $\psi:\mathbb R\to\mathbb R$, define
\[
\psi^*(x)=\sup_{t\in\mathbb R}\{tx-\psi(t)\}
 \in(-\infty,\infty].
\]
\end{de}

\begin{prop}
The transform $\psi^*$ is convex and lower semicontinuous, and a finite convex
function on $\mathbb R$ satisfies $\psi^{**}=\psi$. If
$Y=X-\mathbb EX$, then $\psi_Y(0)=\psi_Y'(0)=0$ and $\psi_Y\geq0$. Hence, for
$x\geq0$,
\[
\psi_Y^*(x)=\sup_{t\geq0}\{tx-\psi_Y(t)\}\geq0.
\]
\end{prop}

\begin{proof}
A supremum of affine functions is convex and lower semicontinuous. The
biconjugacy statement is the supporting-line theorem for finite convex
functions. Finally,
\[
\psi_Y''(t)
 =\operatorname{Var}_{Q_t}(Y)\geq0,
\qquad
\frac{\dd Q_t}{\dd P}=e^{tY-\psi_Y(t)},
\]
so $0$ is a global minimum of $\psi_Y$. For $t\leq0$ and $x\geq0$,
$tx-\psi_Y(t)\leq0$, while $t=0$ gives zero.
\end{proof}

\begin{thm}[One-sided Cram\'er theorem]
Let $(X_n)$ be i.i.d. exponentially integrable random variables. Put
$Y_n=X_n-\mathbb EX_1$, $S_n=\sum_{k=1}^{n}Y_k$, and
$I=\psi_{Y_1}^*$. Assume $Y_1$ is not almost surely constant. For every
\[
0<\varepsilon<\operatorname*{ess\,sup}Y_1,
\]
we have
\[
\lim_{n\to\infty}\frac1n\log
P\!\left(\frac{S_n}{n}\geq\varepsilon\right)
 =-I(\varepsilon).
\]
If $Y_1=0$ almost surely, the same interpretation holds with
$I(\varepsilon)=\infty$ for $\varepsilon>0$.
\end{thm}

\begin{proof}
For every $t\geq0$, Markov's inequality gives
\[
P(S_n\geq n\varepsilon)
 \leq e^{-tn\varepsilon}\mathbb E e^{tS_n}
 =\exp\{-n(t\varepsilon-\psi_{Y_1}(t))\}.
\]
Taking the infimum over $t$ yields the upper bound
$\limsup n^{-1}\log P(S_n\geq n\varepsilon)\leq-I(\varepsilon)$.

For the lower bound, convexity and exponential integrability imply that there
is $t>0$ with $\psi_{Y_1}'(t)=\varepsilon$. Tilt the common law of $Y_1$ by
\[
\frac{\dd\nu_t}{\dd\mathcal L(Y_1)}(y)
 =e^{ty-\psi_{Y_1}(t)},
\]
and denote probability and expectation under its $n$-fold product by
$Q_t^{(n)}$ and $\mathbb E_t^{(n)}$. Under this product law, the coordinates
remain i.i.d., have mean $\varepsilon$, and have finite positive variance. For
$\delta>0$,
\begin{align*}
P(n\varepsilon\leq S_n\leq n(\varepsilon+\delta))
&=\mathbb E_t^{(n)}\!\left[
 e^{-tS_n+n\psi_{Y_1}(t)}
 \mathbf1_{\{n\varepsilon\leq S_n\leq n(\varepsilon+\delta)\}}
 \right]\\
&\geq
 e^{-n(t\varepsilon-\psi_{Y_1}(t))-nt\delta}
 Q_t^{(n)}(n\varepsilon\leq S_n\leq n(\varepsilon+\delta)).
\end{align*}
The central limit theorem under the tilted product laws shows that the final
probability tends to $1/2$. Since
$t\varepsilon-\psi_{Y_1}(t)=I(\varepsilon)$, taking logarithms gives
\[
\liminf_{n\to\infty}\frac1n\log P(S_n\geq n\varepsilon)
 \geq-I(\varepsilon)-t\delta.
\]
Let $\delta\downarrow0$.
\end{proof}
